\documentclass{article}

\usepackage{graphicx}
\usepackage{amsmath}
\usepackage{booktabs}
\usepackage{hyperref}

\title{DNNKit: A Lightweight Framework for Reproducible Deep Learning Experiments}

\author{
Festus M. Slade \\
Cambridge Neuroscience, University of Cambridge, Cambridge CB2 1TN, UK \\
\texttt{github.com/Festus0/dnnkit}
}

\begin{document}

\maketitle

\begin{abstract}
Reproducibility and experiment management remain major challenges in modern deep learning research.
While powerful frameworks such as PyTorch Lightning and HuggingFace Trainer provide extensive tooling,
their complexity can introduce significant overhead for small-scale research projects, educational settings,
and rapid experimentation workflows.
We present \textbf{DNNKit}, a lightweight framework for reproducible deep learning experimentation that integrates
a command-line training interface, dataset registry, experiment logging, benchmark tracking,
and automatic report generation.
DNNKit aims to provide a minimal, transparent infrastructure that simplifies experiment management
while maintaining compatibility with PyTorch.
We demonstrate the functionality of the framework using a benchmark experiment on the MNIST handwritten
digit classification dataset.
\end{abstract}

\section{Introduction}

Deep learning research often involves large numbers of experiments, parameter configurations,
datasets, and model variants.
Managing these experiments in a reproducible and structured way is a persistent challenge.

Common research workflows frequently rely on ad-hoc scripts, notebooks, and manually maintained logs.
This can make it difficult to reproduce results, compare experiments, or maintain clear experiment provenance.

Several frameworks have been developed to address these issues.
For example, PyTorch Lightning provides a structured training interface,
while HuggingFace Trainer focuses on large-scale model training and benchmarking.
Although these frameworks are powerful, they introduce additional layers of abstraction
that may be unnecessary for smaller research projects or educational environments.

In this work, we introduce \textbf{DNNKit}, a lightweight experimentation framework designed
to simplify experiment management while remaining transparent and easy to modify.
The framework focuses on four core components:

\begin{itemize}
\item Command-line training interface
\item Dataset registry
\item Experiment logging and result storage
\item Automatic leaderboard and report generation
\end{itemize}

DNNKit is designed to provide a minimal infrastructure for reproducible experiments
while remaining fully compatible with PyTorch-based models.

\section{Framework Design}

The design philosophy of DNNKit prioritizes simplicity, transparency, and reproducibility.

\subsection{Training CLI}

The framework provides a command-line interface (CLI) that standardizes experiment execution.
Training runs can be launched using a simple command:

\begin{verbatim}
dnnkit train --dataset mnist --epochs 3
\end{verbatim}

Each training run automatically generates a dedicated experiment directory containing:

\begin{itemize}
\item model checkpoints
\item training metrics
\item configuration parameters
\end{itemize}

\subsection{Dataset Registry}

Datasets are accessed through a registry mechanism that standardizes data loading.
This avoids repeated boilerplate code across experiments and ensures consistent dataset usage.

\subsection{Experiment Logging}

All training runs are automatically logged to structured output directories.
Metrics such as training loss, accuracy, and configuration parameters are stored
in machine-readable formats.

\subsection{Leaderboard Tracking}

DNNKit aggregates experiment results into a global leaderboard,
allowing quick comparison across runs.

\subsection{Automatic Report Generation}

The framework can automatically generate experiment reports summarizing results
and benchmark performance.

\section{Experimental Evaluation}

To demonstrate the functionality of DNNKit,
we evaluated the framework using the MNIST handwritten digit classification dataset.

\subsection{Dataset}

MNIST consists of 70,000 grayscale images of handwritten digits
with a resolution of $28 \times 28$ pixels.

\subsection{Model}

We used a simple fully connected neural network consisting of:

\begin{itemize}
\item Input layer (784 features)
\item Hidden layer with ReLU activation
\item Output layer with 10 classes
\end{itemize}

The model was trained using the Adam optimizer.

\subsection{Training Configuration}

\begin{itemize}
\item Optimizer: Adam
\item Learning rate: $1 \times 10^{-3}$
\item Batch size: 64
\item Epochs: 3
\end{itemize}

\section{Results}

Table~\ref{tab:results} summarizes the benchmark results obtained using DNNKit.

\begin{table}[h]
\centering
\begin{tabular}{ccc}
\toprule
Epoch & Training Loss & Test Accuracy \\
\midrule
1 & 0.3442 & 94.46\% \\
2 & 0.1610 & 95.95\% \\
3 & 0.1110 & 96.97\% \\
\bottomrule
\end{tabular}
\caption{MNIST benchmark results obtained using DNNKit.}
\label{tab:results}
\end{table}

The results demonstrate that the framework supports standard deep learning training workflows
while maintaining a lightweight and transparent architecture.

\section{Discussion}

DNNKit aims to fill the gap between simple research scripts and large training frameworks.

Compared to large frameworks, DNNKit offers several advantages:

\begin{itemize}
\item minimal abstraction overhead
\item easy code inspection
\item fast experiment setup
\item reproducible experiment tracking
\end{itemize}

The framework is particularly well suited for:

\begin{itemize}
\item rapid research prototyping
\item small-scale machine learning experiments
\item educational environments
\end{itemize}

Future development may include additional features such as distributed training support,
extended dataset registries, and experiment visualization tools.

\section{Conclusion}

We introduced DNNKit, a lightweight framework for reproducible deep learning experiments.
The framework integrates experiment management, dataset handling, benchmarking,
and report generation into a simple command-line interface.

Through a benchmark experiment on MNIST, we demonstrated that DNNKit provides
a practical infrastructure for organizing and reproducing machine learning experiments.

The source code is publicly available at:

\begin{center}
\url{https://github.com/Festus0/dnnkit}
\end{center}

\section*{Acknowledgments}

This work was developed as part of independent research into lightweight machine learning
infrastructure for reproducible experimentation.

\bibliographystyle{unsrt}

\begin{thebibliography}{3}

\bibitem{pytorch}
Paszke, A. et al.
PyTorch: An Imperative Style, High-Performance Deep Learning Library.
NeurIPS (2019).

\bibitem{mnist}
LeCun, Y. et al.
Gradient-Based Learning Applied to Document Recognition.
Proceedings of the IEEE (1998).

\bibitem{lightning}
Falcon, W.
PyTorch Lightning.
GitHub Repository.

\end{thebibliography}

\end{document}
